\documentclass[11pt,a4paper]{article}

\usepackage[utf8]{inputenc}
\usepackage[T1]{fontenc}
\usepackage[spanish,es-noquoting]{babel}
\usepackage[a4paper,margin=2.3cm]{geometry}
\usepackage{hyperref}
\usepackage{enumitem}
\usepackage{xcolor}

\hypersetup{
  colorlinks=true,
  linkcolor=blue!50!black,
  urlcolor=blue!50!black
}

\setlength{\parindent}{0pt}
\setlength{\parskip}{0.5em}
\setlist[itemize]{leftmargin=1.4em}
\setlist[enumerate]{leftmargin=1.6em}

\title{\textbf{Pr\'actica 04 - SubSonic Festival}\\Dockerizaci\'on y despliegue en Azure}
\author{Programaci\'on en Internet 2025/26}
\date{Abril de 2026}

\begin{document}

\maketitle

\textbf{Integrantes del grupo}

\begin{itemize}
  \item Dar\'io Mateos Mart\'in
  \item Nicol\'as Benito Benito
  \item Manuel Cordero Jaramillo
  \item Gonzalo Naranjo Carretero
  \item Samantha Casta\~neda Godinez
  \item Junior Mart\'inez N\'u\~nez
\end{itemize}

\newpage
\tableofcontents
\newpage

\section{Introducci\'on}

El presente informe recoge el trabajo realizado para la Pr\'actica 04 del caso de estudio SubSonic Festival. Tras haber desarrollado previamente el front-end y el back-end de la aplicaci\'on, esta entrega se centra en preparar el proyecto para su ejecuci\'on mediante contenedores Docker y para su posterior despliegue en un proveedor cloud, concretamente Microsoft Azure.

La aplicaci\'on parte de un proyecto web con front-end en HTML, CSS y JavaScript puro, y back-end en Python con FastAPI. En la versi\'on actual, FastAPI sirve tanto la API REST como los recursos est\'aticos del front-end, por lo que el sistema se puede empaquetar en un \'unico contenedor Docker sin introducir frameworks nuevos ni modificar la arquitectura funcional del proyecto.

El objetivo de esta pr\'actica es que el profesor pueda clonar el repositorio, construir la imagen Docker sin errores, ejecutar la aplicaci\'on en localhost y, posteriormente, desplegarla en Azure siguiendo una gu\'ia reproducible.

\section{Estado de partida del proyecto}

Antes de realizar la dockerizaci\'on, el proyecto SubSonic Festival ya inclu\'ia:

\begin{itemize}
  \item Back-end en Python con FastAPI.
  \item Front-end en HTML, CSS y JavaScript puro.
  \item Persistencia local mediante archivos JSON.
  \item Arquitectura organizada en controller, model, DAO, DTO y factory.
  \item Funcionalidades de Cliente, Proveedor y Visitante integradas con el back-end.
  \item Autenticaci\'on mediante token de sesi\'on tipo bearer para las rutas protegidas.
  \item Prueba autom\'atica general mediante el archivo \texttt{smoke\_test.py}.
\end{itemize}

La dockerizaci\'on se ha realizado respetando esta estructura, sin reescribir el proyecto desde cero y sin separar el front-end en otro servidor.

\section{Gu\'ia del proceso}

\subsection{Estructura Docker incorporada al repositorio}

Siguiendo el enunciado de la pr\'actica, se ha a\~nadido al repositorio una carpeta denominada \texttt{Docker}. Esta carpeta contiene los archivos necesarios para construir y preparar el despliegue de la aplicaci\'on:

\begin{itemize}
  \item \texttt{Docker/Dockerfile}: define la imagen Docker de la aplicaci\'on.
  \item \texttt{Docker/README.md}: resume los comandos de construcci\'on, ejecuci\'on y prueba local.
  \item \texttt{Docker/azure-deploy-template.ps1}: incluye una plantilla de comandos para un despliegue mediante Azure Container Registry.
  \item \texttt{.github/workflows/docker-ghcr.yml}: automatiza la construcci\'on de la imagen Docker y su publicaci\'on en GitHub Container Registry.
\end{itemize}

Adem\'as, se ha a\~nadido un archivo \texttt{.dockerignore} en la ra\'iz del proyecto. Su finalidad es evitar que se copien al contenedor archivos innecesarios como \texttt{.git}, \texttt{.venv}, documentaci\'on generada, archivos temporales o dependencias locales.

\subsection{Imagenes externas utilizadas}

La imagen base utilizada en el Dockerfile es:

\begin{verbatim}
python:3.11-slim
\end{verbatim}

Esta imagen procede de Docker Hub y se utiliza como base ligera para ejecutar aplicaciones Python. Sobre ella se instalan las dependencias definidas en \texttt{requirements.txt} y se copia el c\'odigo necesario para ejecutar SubSonic Festival.

La finalidad de esta imagen externa es proporcionar un entorno Python estable y reproducible sin depender de la configuraci\'on concreta del ordenador del desarrollador.

Adem\'as, para el despliegue final se ha publicado la imagen de la aplicaci\'on en GitHub Container Registry, ya que la suscripci\'on de Azure utilizada bloqueaba la creaci\'on de Azure Container Registry por pol\'itica de la cuenta. La imagen final utilizada por Azure Container Apps es:

\begin{verbatim}
ghcr.io/subsonic-grupo4/l4g1/subsonic-festival:practica04
\end{verbatim}

Su finalidad es contener la aplicaci\'on completa SubSonic Festival, incluyendo back-end FastAPI y recursos est\'aticos del front-end.

\subsection{Contenido principal del Dockerfile}

El Dockerfile realiza las siguientes acciones:

\begin{enumerate}
  \item Usa \texttt{python:3.11-slim} como imagen base.
  \item Define \texttt{/app} como directorio de trabajo.
  \item Copia \texttt{requirements.txt}.
  \item Instala las dependencias del back-end.
  \item Copia las carpetas necesarias del proyecto: \texttt{backend}, \texttt{pages}, \texttt{js} y \texttt{css}.
  \item Copia \texttt{smoke\_test.py} para permitir la validaci\'on dentro del contenedor.
  \item Expone el puerto 8000.
  \item Arranca la aplicaci\'on con Uvicorn escuchando en \texttt{0.0.0.0}.
\end{enumerate}

El comando de arranque usado dentro del contenedor es:

\begin{verbatim}
python -m uvicorn backend.controller.controller:app --host 0.0.0.0 --port 8000
\end{verbatim}

Se utiliza \texttt{0.0.0.0} porque dentro de Docker la aplicaci\'on debe aceptar conexiones desde fuera del contenedor.

\subsection{Construcci\'on de la imagen Docker}

Desde la ra\'iz del repositorio, la imagen se construye con el comando:

\begin{verbatim}
docker build -f Docker/Dockerfile -t subsonic-festival:practica04 .
\end{verbatim}

Este comando utiliza el archivo \texttt{Docker/Dockerfile}, instala las dependencias de Python y prepara el contenedor para ejecutar la aplicaci\'on completa.

\subsection{Prueba local del contenedor}

Antes de desplegar en Azure, se comprueba el funcionamiento del contenedor en local mediante:

\begin{verbatim}
docker run --rm -p 8000:8000 subsonic-festival:practica04
\end{verbatim}

Una vez arrancado, se verifican las siguientes rutas:

\begin{itemize}
  \item \url{http://127.0.0.1:8000/pages/visitante.html}
  \item \url{http://127.0.0.1:8000/pages/login.html}
  \item \url{http://127.0.0.1:8000/docs}
  \item \url{http://127.0.0.1:8000/events}
\end{itemize}

La prueba local confirma que el contenedor sirve correctamente las p\'aginas del front-end, la documentaci\'on Swagger de FastAPI y los end-points principales de la API.

\subsection{Prueba autom\'atica dentro del contenedor}

Tambi\'en se valida la imagen ejecutando el test general del proyecto dentro del propio contenedor:

\begin{verbatim}
docker run --rm subsonic-festival:practica04 python smoke_test.py
\end{verbatim}

El resultado esperado es:

\begin{verbatim}
SMOKE TEST OK
\end{verbatim}

Esta prueba comprueba el acceso p\'ublico a eventos, el login de Cliente y Proveedor, las rutas protegidas, el control por rol y el funcionamiento b\'asico de la API.

\section{Despliegue en Microsoft Azure}

\subsection{Servicio cloud elegido}

Para el despliegue cloud se ha utilizado:

\begin{itemize}
  \item GitHub Container Registry, para almacenar la imagen Docker construida autom\'aticamente desde GitHub Actions.
  \item Azure Container Apps, para ejecutar la aplicaci\'on en Azure y obtener una URL p\'ublica.
\end{itemize}

Esta opci\'on permite desplegar una aplicaci\'on basada en contenedores sin tener que administrar manualmente una m\'aquina virtual. Durante el proceso se intent\'o utilizar Azure Container Registry, pero la suscripci\'on de Azure for Students bloque\'o la creaci\'on de dicho recurso mediante una pol\'itica de regiones/recursos. Por ello, se opt\'o por publicar la imagen en GitHub Container Registry y desplegarla posteriormente en Azure Container Apps como imagen de registro externo.

\subsection{Flujo de despliegue}

El flujo seguido para publicar la aplicaci\'on es:

\begin{enumerate}
  \item Crear la rama \texttt{practica04-docker-azure} para aislar los cambios de la pr\'actica.
  \item A\~nadir el \texttt{Dockerfile}, el \texttt{.dockerignore} y la documentaci\'on asociada.
  \item A\~nadir el workflow \texttt{.github/workflows/docker-ghcr.yml}.
  \item Ejecutar GitHub Actions para construir la imagen Docker.
  \item Publicar la imagen en GitHub Container Registry.
  \item Crear una Azure Container App desde el portal de Azure.
  \item Configurar la Container App para leer la imagen desde \texttt{ghcr.io}.
  \item Activar la entrada externa HTTP en el puerto 8000.
  \item Obtener la URL p\'ublica de la aplicaci\'on desplegada.
\end{enumerate}

Como apoyo, tambi\'en se conserva una plantilla de despliegue mediante Azure Container Registry en:

\begin{verbatim}
Docker/azure-deploy-template.ps1
\end{verbatim}

\subsection{Publicaci\'on de la imagen con GitHub Actions}

La imagen Docker se construye y publica mediante el workflow:

\begin{verbatim}
.github/workflows/docker-ghcr.yml
\end{verbatim}

El workflow se ejecuta al subir cambios a la rama \texttt{practica04-docker-azure} y realiza los siguientes pasos:

\begin{enumerate}
  \item Descarga el repositorio.
  \item Inicia sesi\'on en GitHub Container Registry usando \texttt{GITHUB\_TOKEN}.
  \item Construye la imagen usando \texttt{Docker/Dockerfile}.
  \item Publica la imagen con la etiqueta \texttt{practica04}.
\end{enumerate}

El resultado de GitHub Actions fue correcto, quedando publicada la imagen:

\begin{verbatim}
ghcr.io/subsonic-grupo4/l4g1/subsonic-festival:practica04
\end{verbatim}

\subsection{Configuraci\'on realizada en Azure Portal}

En Azure Portal se cre\'o una aplicaci\'on de tipo Azure Container App con la siguiente configuraci\'on:

\begin{itemize}
  \item Grupo de recursos: \texttt{l4g1-subsonic-rg}.
  \item Nombre de la aplicaci\'on: \texttt{subsonic-app}.
  \item Regi\'on: Spain Central.
  \item Origen de imagen: Docker Hub u otros registros.
  \item Servidor de registro: \texttt{ghcr.io}.
  \item Imagen: \texttt{ghcr.io/subsonic-grupo4/l4g1/subsonic-festival:practica04}.
  \item Tipo de imagen: privada, usando un token de GitHub con permiso de lectura de paquetes.
  \item Entrada externa: habilitada.
  \item Tipo de entrada: HTTP.
  \item Puerto de destino: 8000.
\end{itemize}

\subsection{URL final de despliegue}

Tras crear la aplicaci\'on en Azure, la URL p\'ublica se puede consultar desde el portal o mediante:

\begin{verbatim}
az containerapp show `
  --name subsonic-app `
  --resource-group l4g1-subsonic-rg `
  --query properties.configuration.ingress.fqdn `
  -o tsv
\end{verbatim}

\textbf{URL final de la aplicaci\'on desplegada:} \url{https://subsonic-app.gentlefield-85aec19e.spaincentral.azurecontainerapps.io}

Rutas principales para la comprobaci\'on:

\begin{itemize}
  \item \url{https://subsonic-app.gentlefield-85aec19e.spaincentral.azurecontainerapps.io/pages/visitante.html}
  \item \url{https://subsonic-app.gentlefield-85aec19e.spaincentral.azurecontainerapps.io/pages/login.html}
  \item \url{https://subsonic-app.gentlefield-85aec19e.spaincentral.azurecontainerapps.io/docs}
\end{itemize}

\section{Comprobaciones realizadas}

Durante la preparaci\'on de esta entrega se han realizado las siguientes comprobaciones:

\begin{itemize}
  \item Ejecuci\'on del test general del proyecto en local.
  \item Construcci\'on correcta de la imagen Docker.
  \item Arranque del contenedor en localhost.
  \item Comprobaci\'on de rutas web y API principales.
  \item Ejecuci\'on de \texttt{smoke\_test.py} dentro del contenedor.
  \item Ejecuci\'on correcta del workflow de GitHub Actions.
  \item Publicaci\'on de la imagen en GitHub Container Registry.
  \item Despliegue correcto de la aplicaci\'on en Azure Container Apps.
  \item Comprobaci\'on de la URL p\'ublica final en navegador.
\end{itemize}

El resultado de la prueba autom\'atica ha sido:

\begin{verbatim}
SMOKE TEST OK
\end{verbatim}

\section{Notas para la entrega}

El repositorio debe contener la carpeta \texttt{Docker} y el resto de archivos necesarios para construir la imagen. Adem\'as, seg\'un el enunciado de la pr\'actica, el nombre del repositorio debe seguir la nomenclatura del grupo indicada por la asignatura y se debe invitar al profesor al repositorio de GitHub.

La entrega final debe incluir este informe en formato PDF, el enlace al repositorio donde est\'en subidos los archivos de Docker y la URL p\'ublica de la aplicaci\'on desplegada en Azure.

\section{Anexo: actas de reuni\'on}

\subsection*{Acta 1}

\textbf{Fecha y hora:} 17 de abril de 2026, en clase.

\textbf{Asistentes:} Dar\'io Mateos Mart\'in, Nicol\'as Benito Benito, Gonzalo Naranjo Carretero, Junior Mart\'inez N\'u\~nez y Samantha Casta\~neda Godinez.

\textbf{Contenido:} Revisi\'on en clase del enunciado de la Pr\'actica 04 y an\'alisis de los requisitos de contenedores, Dockerfile, Azure e informe de entrega.

\textbf{Decisiones:} Dockerizar la aplicaci\'on en un \'unico contenedor, ya que FastAPI sirve tanto el back-end como los recursos del front-end.

\subsection*{Acta 2}

\textbf{Fecha y hora:} 24 de abril de 2026, 12:00.

\textbf{Asistentes:} Nicol\'as Benito Benito.

\textbf{Contenido:} Preparaci\'on de la carpeta \texttt{Docker}, creaci\'on del \texttt{Dockerfile}, definici\'on del \texttt{.dockerignore} y validaci\'on local de la imagen.

\textbf{Decisiones:} Utilizar la imagen base \texttt{python:3.11-slim}, exponer el puerto 8000 y ejecutar la aplicaci\'on mediante Uvicorn.

\subsection*{Acta 3}

\textbf{Fecha y hora:} 27 de abril de 2026, 15:30.

\textbf{Asistentes:} Nicol\'as Benito Benito.

\textbf{Contenido:} Revisi\'on del proceso de despliegue en Azure, comprobaci\'on de GitHub Actions, publicaci\'on de la imagen Docker en GitHub Container Registry y cierre de la documentaci\'on de entrega.

\textbf{Decisiones:} Designar a Nicol\'as Benito Benito como miembro encargado de realizar la entrega final de la pr\'actica y de invitar al profesor al repositorio de GitHub.

\end{document}
